% Part: first-order-logic
% Chapter: natural-deduction
% Section: soundness-identity

\documentclass[../../../include/open-logic-section]{subfiles}

\begin{document}

\olfileid{fol}{ntd}{sid}

\olsection{\usetoken{S}{identity} સાથે યથાર્થતા}

\begin{prop}
$\eq$ માટેના નિયમો ધરાવતું સ્વાભાવિક નિગમન યથાર્થ છે.
\end{prop}

\begin{proof}
સ્વરૂપ $\eq[t][t]$નું કોઈ પણ !!{formula} પ્રામાણ્ય છે, કારણ કે દરેક
!!{structure}~$\Struct M$ માટે $\Sat{M}{\eq[t][t]}$. (નોંધો કે પદ $t$
બંધ છે, એટલે કે તેમાં કોઈ ચલ નથી, એમ આપણે માનીએ છીએ; તેથી
ચલ-નિયુક્તિઓ અપ્રસ્તુત છે.)

માનો કે કોઈ !!a{derivation}માં છેલ્લું અનુમાન \Elim{\eq} છે, એટલે કે
!!{derivation}નું સ્વરૂપ નીચે પ્રમાણે છે:
\begin{prooftree}
  \AxiomC{$\Gamma_1$}
  \RightLabel{$\delta_1$}
  \DeduceC{$\eq[t_1][t_2]$}
  \AxiomC{$\Gamma_2$}
  \RightLabel{$\delta_2$}
  \DeduceC{$!A(t_1)$}
  \RightLabel{\Elim{\eq}}
  \BinaryInfC{$!A(t_2)$}
\end{prooftree}
આધારવાક્યો $\eq[t_1][t_2]$ અને $!A(t_1)$ અનુક્રમે !!{undischarged}
ધારણાઓ~$\Gamma_1$ અને $\Gamma_2$માંથી !!{derive}d થયાં છે. આપણે
$!A(t_2)$ એ $\Gamma_1 \cup \Gamma_2$માંથી ફલિત થાય છે તેમ બતાવવું છે.
$\Sat{M}{\Gamma_1 \cup \Gamma_2}$ હોય એવી કોઈ !!a{structure}~$\Struct{M}$
લો. આગમન-પરિકલ્પના મુજબ $\Sat{M}{!A(t_1)}$ અને
$\Sat{M}{\eq[t_1][t_2]}$. તેથી $\Value{t_1}{M} = \Value{t_2}{M}$.
$s$ કોઈ પણ ચલ-નિયુક્તિ હોય અને
$m = \Value{t_1}{M} = \Value{t_2}{M}$ હોય. ત્યારે
\olref[fol][syn][ext]{prop:ext-formulas} મુજબ $\Sat{M}{!A(t_1)}[s]$ જો
અને માત્ર જો $\Sat{M}{!A(x)}[\Subst{s}{m}{x}]$, અને તે જો અને માત્ર જો
$\Sat{M}{!A(t_2)}[s]$. $\Sat{M}{!A(t_1)}$ હોવાથી
$\Sat{M}{!A(t_2)}$ મળે છે.
\end{proof}

\end{document}
